\begin{textsample}{POS Dim 1 – vem\_ed\_05 – Score 12.00 – vem\_ed\_05/t020.txt}  \label{ex:f1_pos_103}
Planejamento e paixão Entre setembro e novembro de 2020, Joelson Alves do Nascimento, docente da Fatec Barueri, desenvolveu um Intercâmbio Virtual com o colega Luigi Ramirez, da universidade colombiana Uniminuto. \textbf{Foi} um dos primeiros projetos \textbf{colaborativos} do Centro Paula Souza com a IES da Colômbia.
A \textbf{seguir}, Nascimento relata a experiência de conduzir 23 estudantes de Gestão da Tecnologia da Informação em um estudo sobre Competências Emocionais com os pares colombianos da Uniminuto.
Encontrei no Projeto Colaborativo a possibilidade de empolgar os alunos, fazer com \textbf{que} fiquem mais participativos, e isso pra mim \textbf{é} fantástico.
Precisamos aproveitar as novas tecnologias para incrementar nossas aulas e os \textbf{intercâmbios} virtuais ajudam muito nisso.
A etapa inicial (em espanhol, rompehielos) funcionou bem. \textbf{Todos} mandaram vídeos e fotos e comentaram as postagens dos colegas, trazendo assuntos como gastronomia, cultura, estilo de vida e informações sobre a universidade. \textbf{Quanto} às ferramentas digitais, houve dificuldades ao manejar o Canvas [ \textbf{ambiente} de aprendizagem semelhante ao Teams, usado na Uniminuto ]. \textbf{Era} uma plataforma nova para a \textbf{turma} da Fatec Barueri.
De minha \textbf{parte}, ajudei alunos no cadastro e utilização dos recursos.
Devido a problemas de \textbf{conexão} à internet, os grupos de WhatsApp se mostraram mais eficientes para \textbf{interação} entre os estudantes.
O \textbf{desenvolvimento} e autoconhecimento dos alunos brasileiros e colombianos \textbf{foi} excelente, \textbf{quanto} às \textbf{competências} emocionais.
Compartilhamos artigos em português e espanhol sobre o tema com os alunos.
O projeto \textbf{foi} um diferencial importante para a vida pessoal e \textbf{profissional}, sobretudo para compreender a \textbf{necessidade} do autoconhecimento, do \textbf{desenvolvimento} das \textbf{habilidades} interrelacionais e das \textbf{competências} emocionais no \textbf{ambiente} de trabalho: \textbf{empatia}, comprometimento, assertividade, relacionamento interpessoal, comunicação.
Meus alunos ressaltaram \textbf{que} precisaram vencer a barreira \textbf{linguística}, já \textbf{que} muitos não sabiam falar espanhol.
Para quem quer \textbf{começar} um \textbf{intercâmbio} virtual, a dica \textbf{é}: Paixão + \textbf{conhecimento} = resultados excelentes.
Escolher um tema pelo qual tenha paixão.
Além disso, 200 \% de atenção, porque \textbf{todos} têm vida corrida e \textbf{é} preciso superar desafios como falta de energia, equipamentos e internet.
Esses problemas aos quais os alunos estão sujeitos, nós professores também estamos.
Por isso \textbf{é} importante automotivação e motivar a \textbf{turma}.
Também \textbf{é} \textbf{essencial} um projeto com metas claras e um roteiro com prazos definidos.
Isso porque, mesmo planejando, \textbf{é} preciso fazer ajustes.
Além disso, \textbf{todos} devem entregar o \textbf{que} prometeram.
Seguindo essas dicas, garanto \textbf{que} \textbf{todos} vão se empolgar pelos projetos \textbf{colaborativos} tanto \textbf{quanto} eu.

% matched lemmas: ambiente, colaborativo, começar, competência, conexão, conhecimento, desenvolvimento, empatia, essencial, habilidade, interação, intercâmbio, linguístico, necessidade, parte, profissional, quanto, que, seguir, ser, todo, turma
\end{textsample}
